
\documentclass[11pt,a4paper]{article}

% Cargar configuración desde archivo externo
% =========================
% PAQUETES
% =========================
\usepackage[utf8]{inputenc}
\usepackage[T1]{fontenc}
\usepackage[spanish]{babel}
\usepackage{geometry}
\usepackage{graphicx}
\usepackage{xcolor}
\usepackage{tikz}
\usepackage{tabularx}
\usepackage{enumitem}
\usepackage{hyperref}
\usepackage{fancyhdr}
\usepackage{parskip}
\usepackage{ragged2e}
\usepackage{helvet}
\usepackage{titlesec}  % <-- Agregar este paquete
\renewcommand{\familydefault}{\sfdefault}

% Justificacion uniforme y mejor ajuste de lineas para evitar desbordes.
\AtBeginDocument{\justifying}
\setlength{\emergencystretch}{1.5em}

% =========================
% GEOMETRÍA
% =========================
\geometry{
    top=2cm,
    bottom=2.5cm,
    left=2.5cm,
    right=2.5cm
}

% =========================
% COLORES PERSONALIZADOS
% =========================
\definecolor{ieeeOrange}{HTML}{E07830}
\definecolor{ieeeGreen}{HTML}{8FBC8F}
\definecolor{ieeeLightBlue}{HTML}{87CEEB}
\definecolor{ieeeGray}{HTML}{C0C0C0}
\definecolor{titleBrown}{HTML}{5D4E37}
\definecolor{sectionTeal}{HTML}{008080}
\definecolor{linkTeal}{HTML}{007070}
\definecolor{textBrown}{HTML}{716650}

% =========================
% FORMATEO DE SECCIONES
% =========================
\titleformat{\section}
    {\titlerule
    \Large\bfseries\color{ieeeOrange}}  % Formato: grande, negrita, naranja
    {}                                    % Sin número
    {0pt}                                 % Sin espacio entre número y título
    {}                                    % Sin prefijo
\titlespacing*{\section}{0pt}{0.8cm}{0.4cm}  % Espaciado: izq, antes, después

\titleformat{\subsection}
    {\large\bfseries\color{ieeeOrange}}
    {}
    {0pt}
    {}
\titlespacing*{\subsection}{0pt}{0.6cm}{0.3cm}

% =========================
% COMANDOS PERSONALIZADOS
% =========================

% Barra de colores superior
\newcommand{\colorbar}{%
    \noindent
    \begin{tikzpicture}[remember picture, overlay]
        \fill[ieeeOrange] (current page.north west) rectangle ++(5.25cm,-0.4cm);
        \fill[ieeeGreen] ([xshift=5.25cm]current page.north west) rectangle ++(5.25cm,-0.4cm);
        \fill[ieeeLightBlue] ([xshift=10.5cm]current page.north west) rectangle ++(5.25cm,-0.4cm);
        \fill[ieeeGray] ([xshift=15.75cm]current page.north west) rectangle ++(5.25cm,-0.4cm);
    \end{tikzpicture}
}

% Línea de acción (color naranja)
\newcommand{\lineaAccion}[2]{%
    \subsection*{Línea de Acción #1: #2}
}

% Fila de tabla descriptiva
\newcommand{\filaDescriptiva}[2]{%
    \textbf{#1} & #2 \\[0.4cm]
}

\setlist[enumerate,1]{label=\Roman*., font=\bfseries, leftmargin=1.5cm}
\setlist[enumerate,2]{label=\roman*., leftmargin=1cm}

\newcommand{\logos}{
% Logos (placeholder - reemplazar con logos reales)
\noindent
\begin{minipage}[c]{0.2\textwidth}
    \centering
    % Logo IEEE (diamante)
    \includegraphics[width=0.9\linewidth]{src/IEEE.png}
\end{minipage}%
\hfill
\begin{minipage}[c]{0.4\textwidth}
    \centering
    % Logo IEEE Aguascalientes Section
    \includegraphics[width=0.9\linewidth]{src/IEEE_AGS.png}
\end{minipage}%
\hfill
\begin{minipage}[c]{0.25\textwidth}
    \centering
    % Logo Young Professionals
    \includegraphics[width=1.2\linewidth]{src/UAA2.png}

\end{minipage}



}

% =========================
% DOCUMENTO
% =========================
\begin{document}
\color{textBrown}

% ============================================================
% PORTADA
% ============================================================
\colorbar
\pagestyle{empty}

\vspace{1.5cm}

\logos
\vspace{3cm}

% Título principal
{\fontsize{36}{42}\selectfont\bfseries\color{titleBrown} Propuesta de evento:\\ Conferencia de difusión}

\vspace{0.8cm}

% Subtítulo
{\large Rama Estudiantil UAA Campus Central}

\vspace{0.3cm}

% Línea decorativa
\noindent\rule{7.5cm}{0.5pt}

\vfill

% Información inferior
{\color{sectionTeal}\bfseries IEEE Sección Aguascalientes}\\[0.2cm]
Última fecha de revisión: 22 de Marzo 2026\\[0.1cm]
Fecha de autorización:

\newpage

% ============================================================
% DATOS DE IDENTIFICACIÓN
% ============================================================
\pagestyle{fancy}
\fancyhf{}
\renewcommand{\headrulewidth}{0pt}
\fancyfoot[C]{\thepage}

\vspace{0.5cm}

\section{Datos de identificación del documento}

\vspace{0.8cm}

\hyphenpenalty=10000
\exhyphenpenalty=10000

\begin{center}
\begin{tabular}{@{}p{4.5cm} p{8cm}@{}}
    \small Categoría & \small Detalle \\[0.5cm]
    
    \textbf{Título del Plan} & Conferencia de difusión de la Rama Estudiantil y Convenio 2026 \\[0.5cm]
    
    \textbf{Sección IEEE} & Aguascalientes (Región 9) \\[0.5cm]
    
    \textbf{Coordinación Responsable} & Rama estudiantil IEEE campus UAA \\[0.5cm]
    
    \textbf{Líder del Proyecto} & Sara Fernández Muñoz\\[0.5cm]
    
    \textbf{Fecha de Redacción} & Marzo de 2026 \\[0.5cm]
    
\end{tabular}
\end{center}


% ============================================================
% CONTENIDO
% ============================================================
\pagestyle{fancy}
\fancyhf{}
\renewcommand{\headrulewidth}{0pt}
\fancyfoot[C]{\thepage}

\vspace{0.5cm}

\section{Introducción y alcance estratégico}

El presente documento formaliza la propuesta de difusión de la Rama Estudiantil IEEE Campus UAA ante jefaturas de departamento y organizaciones estudiantiles consolidadas. Su propósito es establecer, con claridad institucional, el enfoque de trabajo para fortalecer la comunidad académica y profesional de la Universidad Autónoma de Aguascalientes.

La \textbf{Rama Estudiantil IEEE Campus UAA} forma parte de la Sección IEEE Aguascalientes (Región 9) y tiene como misión promover el desarrollo académico, técnico y social de estudiantes de ingeniería y licenciatura mediante actividades de vinculación, formación y colaboración.

\textbf{Alcance del documento.} Esta propuesta delimita objetivos y criterios de coordinación; no sustituye los programas académicos formales ni las funciones administrativas de cada departamento.




% \textbf{Alcance Estratégico (Replicabilidad y Colaboración):}

% % Esta parte la hizo el das 
% Este plan no solo define las actividades operativas de 2026, sino que también establece una estructura metodológica y de gestión estandarizada para ser fácilmente replicada y transferida a futuros liderazgos. Además, la ejecución de las acciones propuestas se llevará a cabo en estrecha colaboración con otras ramas estudiantiles, profesionales y la industria local, buscando sinergias que potencien el impacto de cada iniciativa y fortalezcan el ecosistema IEEE en Aguascalientes.


\section{Objetivo principal}

Consolidar una \textbf{comunidad estudiantil sólida, activa y colaborativa} que impulse el intercambio de conocimiento, la vinculación con profesionales y el desarrollo de competencias técnicas y de liderazgo en la Universidad Autónoma de Aguascalientes.

La Rama Estudiantil IEEE Campus UAA se plantea como una plataforma de apoyo y articulación institucional: complementa esfuerzos ya existentes, promueve la participación estudiantil organizada y fortalece el ecosistema académico sin duplicar funciones de otras sociedades u organizaciones.




\textbf{Mesa Directiva y roles}

La Rama cuenta con una mesa directiva establecida y en funciones. No obstante, el modelo de trabajo propuesto prioriza la participación activa de sus integrantes en la planeación, ejecución y seguimiento de actividades, con enfoque en liderazgo compartido, responsabilidad operativa y trabajo colaborativo.

Asimismo, se contempla la incorporación de roles de apoyo conforme a las necesidades del plan anual, con el fin de fortalecer la gestión interna y diversificar la toma de decisiones.




\section{Justificación estratégica}
Este Plan de Trabajo se fundamenta en dos pilares para el crecimiento sostenido de la Sección IEEE Aguascalientes: el desarrollo integral del capital humano y la sostenibilidad institucional.
\begin{enumerate}
    \item \textbf{Visión humana (bienestar como motor de la excelencia técnica).} La ingeniería es el eje de trabajo, pero su impacto depende de una comunidad estudiantil integrada y participativa. Fortalecer espacios de convivencia, networking y colaboración mejora la motivación, la permanencia y el compromiso con los estándares de excelencia que promueve IEEE.
    \item \textbf{Sostenibilidad y transferencia de conocimiento.} La propuesta adopta una estructura replicable para institucionalizar prácticas de planeación y seguimiento. Esto facilita la continuidad entre generaciones de liderazgo, reduce la curva de aprendizaje operativa y protege la permanencia de iniciativas estratégicas.

\end{enumerate}
\section{Líneas de acción y actividades}

La estrategia operativa se centra en tres líneas de acción fundamentales para el fortalecimiento de la comunidad.

\begin{itemize}[leftmargin=1.5cm]
    \item \textbf{Nota Estratégica:} La ejecución exitosa de estas actividades tiene un doble propósito: la autofinanciación de eventos (ej. Posada) y el aumento de la base de miembros para incrementar la asignación de fondos por parte de la IEEE a la Sección Aguascalientes.
\end{itemize}

\section{Descripción del Evento de Reclutamiento}
Como eje central de la estrategia de expansión, se llevará a cabo una jornada de reclutamiento masivo dirigida a las áreas de ingeniería y tecnología.

\textbf {Sede Propuesta:} Auditorio Pedro de Alba, Campus Central UAA.

\textbf {Público Objetivo:} Estudiantes de las carreras de Ingeniería en Sistemas Computacionales, Ingeniería Electrónica, Ingeniería Industrial, Ingeniería en Ing. en Computación Inteligente, Lic. en Informática y Tecnologías Computacionales.

\textbf {Dinámica de Integración:} El evento consistirá en una presentación de alto impacto sobre los beneficios de la membresía global, el acceso a recursos técnicos y la red de contactos internacional.

\textbf {Estrategia de Retención y Networking}
Para fomentar el sentido de pertenencia desde el primer contacto, el evento concluirá con un espacio de convivencia diseñado para el intercambio de ideas entre aspirantes y miembros actuales:

Aperitivo Institucional: Se ofrecerá café y pan dulce a todos los asistentes, transformando el cierre del evento en un espacio de networking relajado.

Identidad de Marca (Merchandising): Se realizará la entrega de artículos oficiales (merch) de la Rama Estudiantil IEEE a los nuevos interesados, con el fin de proyectar una imagen sólida y atractiva frente a la comunidad universitaria.

\section{Estrategia de Difusión y Publicidad}

Para asegurar la asistencia de las tres carreras convocadas, se implementarán las siguientes acciones:

\textbf {Publicidad Física:} Colocación de flyers oficiales en tableros estratégicos y áreas comunes del campus.

\textbf {Promoción en Aula: } Visitas breves a los salones de clase para invitar directamente a los estudiantes y resolver dudas.

\section{Gestión Administrativa y Peticiones}
Se solicita formalmente a la Jefatura de Departamento el apoyo con los siguientes recursos:

\textbf {Disponibilidad y solicitud del Espacio:} Determinación de fechas y horarios óptimos para el uso del Auditorio Pedro de Alba.

\textbf {Red de Enlaces:} Acceso a la base de datos de tutores y jefes de grupo, así como la identificación de profesores que puedan fungir como vocales y promotores del evento en sus respectivas secciones.

\end{document}
